% =====================================================================
%  GMN3 §4.4--4.5  纯 SU(2) 规范理论中的形式线算子 / 物理 vs. 形式线算子
%  原文：Lamination-Paper.tex 2607--3114 行；印刷页 38--46（PDF p.39--47）
%  编号公式 25 条：原文 2617, 2626, 2637, 2645, 2705, 2716, 2726, 2743,
%    2760, 2778, 2799, 2810, 2820, 2871, 2874, 2887, 2937, 2947, 2960,
%    2976, 2986, 2991, 2998, 3007, 3103 行。
%    （BRIEF 记 26 条：多出的一条是原文 2911 行被 % 注释掉的 eq:8fold-SU2，
%      该段按施工指令不译。）
%  图 3 幅：SU2-FL-weak / SU2-FL-Wilson / SU2-FL-Monomials
%  本块承接 ch05（§4.1--4.3），故不含 \section；式号与图号由 LaTeX 续接。
%  跨块引用一律硬编码原文编号，行尾 % TODO 合稿改 \ref/\eqref。
% =====================================================================

\subsection{纯 $SU(2)$ 规范理论中的形式线算子}
\label{subsec:formal-SU2}

现在转向规范群李代数为 $su(2)$、且不带味的规范理论。乍看之下，它们与 $N=3$ Argyres--Douglas 理论相似得颇具迷惑性。为描述腔结构与 BPS 谱，我们依赖文献 \cite{Gaiotto:2009hg} 的第 10.1 节。
一如既往，$\CB=\IC-\{u_+,u_-\}$，而电荷格的最佳描述是
$\Gamma=H_1(\Sigma_u,\IZ)^-$，其中 $\Sigma_u$ 现在是 $\CB$ 上某一族黎曼面的纤维，该族由
%
\begin{equation}
\lambda^2 = \left(\frac{\Lambda^2}{z^3 } + \frac{2u}{z^2} + \frac{\Lambda^2}{z} \right) (dz)^2.
\end{equation}
%
定义，其中 $z\in \IC^*$。纤维 $\Gamma_u$ 仍是秩二的，向量 $p\gamma_1 + q\gamma_2$ 有时记作 $(p,q)$。与
$N=3$ Argyres--Douglas 理论一个非常重要的差别在于，现在
%
\begin{equation}\label{eq:SU2-innprod}
\langle \gamma_1, \gamma_2 \rangle = 2.
\end{equation}
%
在 $u=u_\pm = \pm \Lambda^2$ 处有两个奇点，$\Gamma$ 绕它们有单值性。（见文献 \cite{Gaiotto:2009hg} 的 (10.4)、(10.5) 两式。）
底空间 $\CB$ 应划分为三个区域 $\CB_s, \CB_{w}^\pm$，在每一区域上 $\Gamma$ 都可平凡化。$\CB_{s}$ 是 $u=0$ 附近的一个区域，$\CB_{w}^\pm$ 是 $u=\infty$ 附近的两个区域。有一道穿过
$u_\pm$ 的连通墙，它把 $\CB_s$ 与 $\CB_{w}^\pm$ 隔开。对 $u \in \CB_s$ 有
%
\begin{equation}
\Omega(\gamma;u) = \begin{cases} 1 & \gamma = \pm \gamma_1, \pm \gamma_2 \\
0 & {\rm else} \\ \end{cases}
\end{equation}
%
内积 \eqref{eq:SU2-innprod}
导致比 $N=3$ 理论更复杂的穿墙：对 $u \in \CB_w^+$，我们发现
%
\begin{equation}\label{eq:SU2-weak-spec}
\Omega(\gamma;u) = \begin{cases} -2 & \gamma = \pm (\gamma_1 +  \gamma_2) \\
1 & \gamma = \pm [\gamma_1 + n(\gamma_1+\gamma_2) ], \quad n \geq 0 \\
1 & \gamma = \pm [\gamma_2 + n(\gamma_1+\gamma_2) ], \quad n \geq 0 \\
0 & {\rm else} \\ \end{cases}
\end{equation}
%
$\CB_w^-$ 中的谱可由单值性或者由穿墙得到。

\insfig{SU2-FL-weak}{纯 $SU(2)$ 理论在弱耦合区域（$u$ 取该区域）的 BPS 墙，投影到 $\zeta$ 平面上。为清晰起见，标记各面墙时省略了 $\widehat{W}$ 上的帽子。在第一象限中，荷为 $n\gamma_1 + (n+1)\gamma_2$
与 $(n+1)\gamma_1 + n\gamma_2$（$n\geq 0$）的墙从两侧向 $\gamma_1 + \gamma_2$ 的墙聚集。除 $\pm (\gamma_1 +\gamma_2)$ 处的中央墙以外，每一面这样的墙上的穿墙都由单个超多重态生成。}

现在可以算出：在 $u \approx 0$ 处存在一组满足 \eqref{eq:SU2-innprod} 的基回路 $\gamma_1, \gamma_2$ 的合适选取，使得 $\zeta$ 平面中的腔结构与图 5% TODO 合稿改 \ref{fig:ADN3-FL-1}
相同。另一方面，在弱耦合下情况截然不同。对固定的 $u \in \CB_w^+$，各腔到 $\zeta$ 平面的投影
见图 \ref{fig:SU2-FL-weak}。可以选一条通向大 $u$ 的路径，把 $\gamma_1, \gamma_2$
与通常的磁荷基、电荷基作比较。由此可以辨认出
$\gamma_1 = -\half H_\alpha +\alpha$ 与 $\gamma_2 =\half H_\alpha$，于是 $\gamma_1+\gamma_2$ 就是
$W$ 玻色子的荷。

设 $u \in \CB_s$。与 $N=3$ 的情形一样，容易检验哪些单项式 $X_{(p,q)}$ 在顺时针与逆时针两个方向的单次穿墙之后仍与强正性相容。结果同样是：

\begin{enumerate}

\item $n=1\mod 4$ 时，$c_n$ 中 $p\geq 0, q\geq 0$

\item $n=2\mod 4$ 时，$c_n$ 中 $p\leq 0, q\geq 0$

\item $n=3\mod 4$ 时，$c_n$ 中 $p\leq 0, q\leq 0$

\item $n=0\mod 4$ 时，$c_n$ 中 $p\geq 0, q\leq 0$

\end{enumerate}

在本节余下部分，我们将证明：在任意穿墙之下，这些单项式都生成 $X_{\gamma}$ 的有限展开。根据第 4.1 节% TODO 合稿改 \ref{subsec:Chambers}
末尾的评注，这在猜想的意义上足以证明强正性。下面我们会为这一假设给出一些合理性论证。利用第 5.2 节% TODO 合稿改 \ref{subsec:Seeds-from-N=2}
末尾的推理\zy{原文 §5.2 与 §5.3 两处共用 label \texttt{subsec:Seeds-from-N=2}，故原文 PDF 把此处的交叉引用印成「第 5.3 节」；按 label 名称与文意应为第 5.2 节，译本取后者，并记入勘误表。}，我们预期：把它们延拓为弱耦合区域中的强正性算子并无障碍。

若我们从某个单项式出发去构造强正性的形式线算子，很快就会发现表达式相当复杂，即使在强耦合区内也是如此。我们将看到，在 $\zeta$ 平面中一般并没有有限的单值性。先从一例开始。方便的记号是
%
\begin{equation}
[a,b]:= X_{\half(a \gamma_1 + b\gamma_2)},
\end{equation}
%
并考虑由 $F(c_1)=[1,0]$ 出发定义的形式线算子。现在沿顺时针方向穿墙，逐个生成其他腔中的值。把各项排成三角阵列会很有启发性：这样一来，就容易看出 halo 如何穿过墙被湮灭。沿顺时针方向，得到 $F(c_2)$ 为：

\begin{equation}\label{eq:SU2-sc-rot1}
\begin{array}{cc}
 &[1,2]\\
 &+ \\
 & [1,0]
\end{array}
\end{equation}
%
$F(c_3)$ 为
%
\begin{equation}\label{eq:SU2-sc-rot2}
 \begin{array}{ccccc}
   [-3,2] & + & \rho_2 [-1,2] &  + & [1,2]\\
    &  & &  & +\\
    & &  & & [1,0]
  \end{array}
 \end{equation}
%
这里为节省书写，我们把一个表示的特征标等同于该表示本身；于是 $\rho_N$ 代表 $[N]$。
当表达式变长时，这很有用。
\eqref{eq:SU2-sc-rot2} 中行与列的各项都可以看成以 $[1,2]$ 为核心电荷的 halo 构型。
接着进入腔 $c_4$，竖直列被消掉，因为一个 halo 被移除；但同时因 halo 的产生而生成另外两个竖直列，于是 $F(c_4)$ 为
%
\begin{equation}
\begin{array}{cccccc}
  &[-3,2] & + & \rho_2 [-1,2] & + & [1,2] \\
     &  + &   & + &   &   \\
 & \rho_3 [-3,0] &   & \rho_2 [-1,0] &  &  \\
  &+ &  &   &   &   \\
  & \rho_3 [-3,-2] &   &   &   &   \\
  &+ &   &  &   &   \\
  &[-3,-4] &   &   &   &
\end{array}
\end{equation}
%
我们再次看到，角上的项 $[-3,2]$ 充当水平 halo 与竖直 halo 的核心粒子。
现在穿过 $\widehat{W(\gamma_1)}$ 进入下一个腔，第一行的水平 halo 塌缩到它的核心，于是得到 $F(c_5)$
%
\begin{equation}
\begin{array}{cccccccccc}
  &[-3,2] &   &   &   &  &   &   &   &   \\
 & + &   &   &   &  &   &   &   &   \\
 & \rho_3 [-3,0] &   & \rho_2 [-1,0] &   &   &   &   &   &   \\
 &  + &   &  + &   &   &   &   &   &   \\
 & \rho_3 [-3,-2] & + & (\rho_2+\rho_4)[-1,-2] & + & \rho_3[1,-2] &   &   &   &   \\
 &  + &   &  + &   &   &   &   &   &   \\
 & [-3,-4] & +& \rho_4 [-1,-4] & + & \Lambda^2 \rho_4 [1,-4] & + & \rho_4 [3,-4] & + & [5,-4]
\end{array}
\end{equation}
%
这个表达式有一个有趣之处：我们必须用到表示之间的恒等式 $\rho_2 \rho_3 = \rho_2 + \rho_4$。为认出
第三行的水平 halo，我们用 $\rho_2 \rho_3$；而为认出
第二列的两个竖直 halo，必须用 $\rho_2 + \rho_4$。作为最后一个例子，
继续沿顺时针穿过 $\widehat{W(\gamma_2)}$，得到 $F(c_6)$ 为
%
\begin{equation}
\begin{array}{ccccccccccc}
    &   &   &   &   &   &   &   &   &   &  [5,6] \\
    &   &   &   &   &   &   &   &   &   &  + \\
    &   &   &   &   &   &   &   &   &   &  \rho_5 [5,4] \\
    &   &   &   &   &   &   &   &   &   &  + \\
    &   &   &   &   &   &   &   &  \rho_4[3,2]  &   &  \Lambda^3\rho_5 [5,2] \\
    &   &   &   &   &   &   &   & +  &   &  + \\
    &   &   &   &   &   &\rho_3 [1,0]   &   & \rho_4\rho_3 [3,0]  &   &  \Lambda^2\rho_5 [5,0] \\
    &   &   &   &   &   &+  &   & +  &   &  + \\
    &   &   &   & \rho_2 [-1,-2]  &   &(\rho_3+\Lambda^2\rho_4) [1,-2]   &   & \rho_4\rho_3 [3,-2]  &   & \rho_5 [5,-2] \\
        &   &   &   &   &   &+  &   & +  &   &  + \\
         &   & [-3,-4]  &  + & \rho_4 [-1,-4]  & +  & \Lambda^2\rho_4 [1,-4]   &  + & \rho_4 [3,-4]  & +  &  [5,-4]
\end{array}
\end{equation}
%
现在完全不清楚的是：穿过墙 $\widehat{W(-\gamma_1)}$ 进入腔 $c_7$ 时，强正性是否仍然保持。有问题的
是倒数第二行。我们必须写 $\rho_3 + \Lambda^2\rho_4 = \rho_2^2 + \rho_5$ 与
$\rho_4 \rho_3 = \rho_5 \rho_2 + \rho_2 $，于是这一行可以写成：
%
\begin{equation}
\rho_2 ([-1,-2] + \rho_2 [1,-2] + [3,-2]) + \rho_5 ([1,-2] + \rho_2 [3,-2] + [5,-2]) .
\end{equation}
%
写成这种形式，就很清楚：水平 halo 将会塌缩，强正性得以保持。

沿逆时针方向推进，我们看到类似的图样。从
$c_1$ 穿过 $\widehat{W(\gamma_1)}$ 到 $c_0$，简单地得到 $F(c_0)=[1,0]$；随后 $F(c_{-1})$ 为
%
\begin{equation}
\begin{array}{cc}
& [1,0] \\
& + \\
& [1,-2]
\end{array}
\end{equation}
%
过渡到 $F(c_{-2})$，得到
%
\begin{equation}
\begin{array}{cccccc}
  &  &   &   &   &  [1,0] \\
  &  &   &   &   & + \\
  & [-3,-2] & + & \rho_2 [-1,-2] & + & [1,-2]
\end{array}
\end{equation}
%
如此等等。

一般的图样如下。从 $n=1 \mod 4$ 的腔 $c_n$ 出发，我们从一个直角三角形开始，
其顶点为 $[a,b], [a-2b,b], [a,b-2a]$，其中 $a,b$ 满足
一些不等式 $a\leq 0, b\leq 0, a-2b \geq 0, b-2a\geq 0, \cdots$。
随后，穿墙在每一个腔中创生并湮灭 halo，形成一个直角三角形阵列，使得
四次穿墙之后我们回到一个三角阵列，其参数为
$a' = - 3a +4 b$、$b' = -4 a + 5b$。在上面的例子中，经 $k$ 次顺时针转动后
$a= 1-4k$、$b= -4k$。直接检验这些三角形内部格点处的强正性
似乎非常不平凡；但只要三角形的图样持续下去，情形必然如此。也就是说，
证明强正性的主要困难在于说明：在穿墙之下，表达式仍保持为
$[a,b]$ 的多项式。事实上，我们将在本节末尾给出一个严格的论证，
说明这种有限性确实成立。

\insfig{SU2-FL-Wilson}{强耦合下的 Wilson 线算子。对这个表达式而言，$\zeta$ 平面中没有单值性。
角上的电荷同时充当水平 halo 与竖直 halo 的核心电荷。
每穿过一面墙，就有一个 halo 被消灭、一个 halo 被创生。}

上述表达式虽已变得相当可观，但存在一个值得注意的特殊情形：在强耦合区，
三角形的参数在 $\zeta$ 平面中没有单值性，
即 $a' = -3a+4b = a $、$b' = -4a + 5 b = b$，
也就是 $a=b$ 且 $a,b<0$。最基本的情形 $a=b=-1$ 将对应于物理理论中基础表示下的 Wilson
线。它在四个腔中的值见图 \ref{fig:SU2-FL-Wilson}。Wilson 算子的
表达式足够简单，以致我们也能在弱耦合区检验强正性。我们从 $c_0$ 与 $c_2$ 中的表达式出发，这两者在
强耦合区与弱耦合区 $\CB_w^+$ 中相同。从 $c_0$ 沿顺时针方向穿过，
在 $n\gamma_1 + (n+1)\gamma_2$ 的墙旁边得到形如 $[1,1]+[2n-1,2n+1] + [-1,-1]$ 的表达式；顺时针穿过这面墙，
一个 halo 被消灭、一个被创生，得到 $[1,1]+[2n+1,2n+3] + [-1,-1]$。
类似地，从 $c_2$ 沿逆时针方向穿过，得到表达式
$[1,1]+[2n+1,2n-1] + [-1,-1]$；再穿过 $(n+1)\gamma_1 + n\gamma_2$ 的墙，
得到 $[1,1]+[2n+3,2n+1] + [-1,-1]$。这样，我们就明确构造出了
对应于 Wilson 圈的形式线算子。

这个 Wilson 线算子（我们记作 $\hat W$）可用来生成一串
形式线算子。我们递归地定义
%
\begin{equation}\label{eq:rec-rels-1}
G_{2k} \hat W = y G_{2k+1} + 1 + y^{-1} G_{2k-1},
\end{equation}
\begin{equation}\label{eq:rec-rels-2}
G_{2k+1}\hat W = y G_{2k+2} + y^{-1} G_{2k}.
\end{equation}
%
递归的出发点是：在腔 $c_1$ 中取 $G_0 = [1,1]$ 的那个算子。（注意，它正是 Wilson
算子的四个顶点所构成的方阵中「缺失」的那个算子。）
回忆 $[a,b]:= X_{\half(a \gamma_1 + b\gamma_2)}$，
于是 $[a,b]\cdot [c,d] = y^{\half(ad-bc)} [a+c,b+d]$。把这些乘法
做出来，就得到一串算子，它们在腔 $c_1$ 中的值为

%
\begin{equation}
\begin{split}
G_{-5}& = [ -2,-4] + \rho_4 [0,-4] + \Lambda^2 \rho_4 [2,-4] + \rho_4 [4,-4] + [ 6,-4] \\
&\quad + \rho_2 ( [-2,-2] + \rho_2 [0,-2]+[2,-2]) + [-2,0]\\
G_{-4}& = [-1,-3] + \rho_3 [1,-3] + \rho_3 [3,-3] + [5,-3] + [-1,-1] + [1,-1] \\
G_{-3} & = [0,-2] + \rho_2 [2,-2] + [4,-2] \\
G_{-2} & = [1,-1] + [3,-1] \\
G_{-1} & = [2,0] \\
G_{0} & = [1,1] \\
G_1 & = [0,2] \\
G_2 & = [-1,1] + [-1,3] \\
G_3 & = [-2,0] + \rho_2 [-2,2] + [-2,4] \\
G_4 & = [-3,-1] + \rho_3 [-3,1] + \rho_3 [-3,3] + [-3,5] + [-1,-1] + [-1,1] \\
G_5 & = [-4,-2] + \rho_4 [-4,0] + \Lambda^2 \rho_4 [-4,2] + \rho_4 [-4,4] + [-4,6] \\
&\quad + \rho_2 ( [-2,-2] + \rho_2 [-2,0]+[-2,2]) + [0,-2]
\end{split}
\end{equation}
%
（再说一次，若把它们排成 $2\times 2$ 的网格，就构成三角阵列。）

\insfig{SU2-FL-Monomials}{本图说明穿墙如何把 $G_n$ 变成单项式。
例如若 $n=8k$，则从 $c_1$ 出发沿顺时针方向做 $4k$ 次穿墙会把 $G_n$
变成单项式 $[1,1]$。}

% 原文 2910--2927 行有一整段公式（label eq:8fold-SU2）被 % 注释掉（未发表内容），按施工指令不译。

这些表达式又一次变得相当复杂。不过它们可以被控制住，
因为足够多次穿墙会把它们带到单项式的形式。设 $\Omega$ 表示
沿顺时针方向穿过一面墙的运算。
那么我们断言：这些穿墙按图 \ref{fig:SU2-FL-Monomials} 所示的图样把 $G_n$ 变成单项式。该图样
可以用归纳法、连同两组各四个非常简单的恒等式来证明，这些恒等式直接验算即可：
%
\begin{equation}\label{eq:4fold-wilson-idents}
\begin{split}
c_{1 + 4\ell} \qquad  & [1,1]\hat W = y[0,2] + 1 + y^{-1} [2,0] \\
c_{2 + 4\ell} \qquad  & [-1,1]\hat W = y[-2,0] + 1 + y^{-1} [0,2] \\
c_{3 + 4\ell} \qquad  & [-1,-1]\hat W = y[0,-2] + 1 + y^{-1} [-2,0] \\
c_{4 + 4\ell} \qquad  & [1,-1]\hat W = y[2,0] + 1 + y^{-1} [0,-2] \\
\end{split}
\end{equation}
%
%
\begin{equation}
\begin{split}
c_{1 + 4\ell} \qquad  & [2,0]\hat W = y[1,1]  + y^{-1} \Omega[1,-1], \\
c_{2 + 4\ell} \qquad  & [0,2]\hat W = y[-1,1]  + y^{-1} \Omega[1,1], \\
c_{3 + 4\ell} \qquad  & [-2,0]\hat W = y[-1,-1]  + y^{-1} \Omega[-1,1], \\
c_{4 + 4\ell} \qquad  & [0,-2]\hat W = y[1,-1]  + y^{-1} \Omega[-1,-1].
\end{split}
\end{equation}
%
由这些恒等式可推出若干高度非平凡的结果。首先，由于 $G_{2n+1}$
在 $\Omega^n$ 之下的穿墙像容许一个平方根，算子本身
也就容许平方根。我们令 $\hat V_n := \sqrt{G_{2n+1}}$。于是例如
%
\begin{equation}
\begin{split}
\hat V_{-3} = \sqrt{G_{-5}} & = [-1,0] + [-1,-2] + \rho_2 [1,-2] + [3,-2],\\
\hat V_{-2} = \sqrt{G_{-3}} & = [0,-1] + [2,-1],\\
\hat V_{-1} = \sqrt{G_{-1}} & = [1,0],\\
\hat V_0 = \sqrt{G_1} & = [0,1],\\
\hat V_1 = \sqrt{G_3} & =    [-1,0] + [-1,2],\\
\hat V_2 = \sqrt{G_5} & =    [0,-1] + [-2,-1] + \rho_2 [-2,1] + [-2,3].
\end{split}
\end{equation}

通过到一个表达式变成单项式（或单项式的简单 $\Omega$ 像）的方便腔，我们可以容易地
建立许多有用的环关系。首先，$G$ 与 $\hat V$ 由下式相联系：
%
\begin{equation}\label{eq:SU2-formring-1}
\begin{split}
G_{2n+1} & = \hat V_n^2 \\
G_{2n} & = y^{-1/2}
\hat V_{n-1} \hat V_n \\
\end{split}
\end{equation}
%
其次，$\hat V$ 满足以下关系：
%
\begin{equation}\label{eq:SU2-formring-2}
\hat V_{n-1} \hat V_{n+1} = 1+ y \hat V_n^2
\end{equation}
%
%
\begin{equation}\label{eq:SU2-formring-3}
  \hat V_{n} \hat W =  y^{1/2} \hat V_{n+1} + y^{-1/2} \hat V_{n-1}
\end{equation}
%
我们稍后会看到，有时用 $G$ 的子代数来工作是有益的。用上面的方法可得，对奇指标的算子有
%
\begin{equation}\label{eq:SU2-formring-4}
\begin{split}
G_{2k-1} G_{2k+1} & = y^2 G_{2k}^2 \\
G_{2k-1} G_{2k+3} & = (1+ y G_{2k+1})(1+y^3 G_{2k+1}) \\
\end{split}
\end{equation}
%
而对偶指标的算子有
%
\begin{equation}\label{eq:SU2-formring-5}
G_{2k} G_{2k+2} = y G_{2k+1} (1+ y G_{2k+1})
\end{equation}
%
如此等等。

在下面的第 8.3 节% TODO 合稿改 \ref{subsec:SU2-Functions}
中，我们将看到这些关系的 $y=1$ 特殊化
如何漂亮地重现 $SU(2)$ 与 $SO(3)$ 规范理论模空间上的函数环。在第 10.3 节% TODO 合稿改 \ref{subsec:SU2-Lams}
中，特别是下面 (10.25) 式% TODO 合稿改 \eqref{eq:chebyshev-expls}
中，我们将从叶层重现
$y=+1$ 时的一串算子 $\hat V_n$。
这样，上面的 $\hat V_n$ 就给出了 $SU(2)$、$N_f=0$ 理论中保护自旋特征标的计算。

现在我们有条件给出一个（在我们看来）有说服力的论证，说明 $\hat V_n$ 与 $G_{2n}$
确实定义了强正性的形式线算子。当强正性失效时，通常会出现
一个无穷级数（例如见 (4.3)、(4.4) 两式）。% TODO 合稿改 \eqref{eq:explicit-dilog-conj}、\eqref{eq:explicit-dilog-conj-ii}
因此，在我们的论证中，我们把有限性当作强正性的替代品。我们希望考虑
$G_n$ 的各种穿墙；特别是，仅从规则 (4.3) 与 (4.4) 出发，并不能明显看出表达式会在某个腔 $c_m$ 中是有限的。% TODO 合稿改 \eqref{eq:explicit-dilog-conj}、\eqref{eq:explicit-dilog-conj-ii}
然而，利用递归关系，我们总可以把 $G_n$ 表成
$G_\ell, G_{\ell+1}, W, y^{\pm 1}$ 的有限多项式，其中 $\ell$ 取得使 $G_{\ell}$ 与 $G_{\ell+1}$
在腔 $c_m$ 中退化为单项式。这就说明 $G_n$ 在每个腔中都是有限的。
对 $\hat V_n$ 有类似的论证。

于是我们得出结论：$G_n$ 与 $\hat V_n$ 全都是强正性的形式线算子。
此外，任何由某个腔中的单项式 $[p,q]$ 生成的线算子，都可以写成
适当的 $\hat V_n$ 的乘积。例如，若我们在腔 $c_n$ 中且 $n=1 \mod 4$，
则 $p,q\geq 0$，并可取 $k$ 使 $[p,q] = y^{-pq/2} \hat V_{4k-1}^p \hat V_{4k}^q$。
由此可知，$[p,q]$ 的所有穿墙像都是 $[a,b]$ 的有限多项式，
其系数在 $su(2)$ 的表示环中。因此，给定上面主要的假设，
$[p,q]$ 确实生成一个强正性的形式线算子。

正如我们前面提到的，这些论证适用于强耦合区中的 $u$。我们
并未试图证明弱耦合区中的强正性，但我们预期它成立，
因为第 8.3 节% TODO 合稿改 \ref{subsec:SU2-Functions}
与第 10.3 节% TODO 合稿改 \ref{subsec:SU2-Lams}
所述与物理线算子的联系。



\subsection{物理线算子 vs. 形式线算子}
\label{subsec:Physics-Formal}

在第 3.2 节中，我们陈述了四个猜想：% TODO 合稿改 \ref{subsec:BPS-Deg-PSC}
关于带框 BPS 态与「原味」BPS 态的
强正性猜想与弱正性猜想。本节我们再陈述一个猜想：
它若成立，就使人能够在已知「原味」BPS 谱的前提下，仅用纯代数的手段
导出全部简单的物理线算子——即使该理论不是拉氏型的。

带框 BPS 态的强正性猜想蕴涵：物理线算子的生成函数 (3.38) 定义一个强正性的形式线算子。% TODO 合稿改 \eqref{eq:form-gen}
我们愿意猜想其逆也成立，即：任何强正性的形式线算子都是
某个物理线算子的带框 BPS 简并度所构成的生成函数。
支持这一猜想的主要证据是：正如我们将在第 10 节% TODO 合稿改 \ref{sec:Examples}
看到的，它在若干例子中成立。

利用单值群的作用，可以对这一猜想作一个漂亮的相容性检验。
物理线算子 $L$ 的具体紫外标号
必须与红外数据 $u$ 无关；只有当沿规范耦合模空间中非平凡的同伦路径搬运时，
它才可能受到单值变换的影响。当在 $\zeta$ 平面中绕原点搬运时，它也可能
受到单值变换的影响，不过只在 $U(1)_R$ 反常的理论中。原因是：
在此类理论中，$\zeta$ 的旋转可以换成
规范耦合能标的旋转，即 theta 角的平移，而这可以引起
紫外标号的单值变换。形式线算子所经历的单值变换，其形式是
单值变换过程中穿过的所有 BPS 墙对应的 KS 变换之乘积。
这与场论中 KS 变换的已知性质相容\cite{Gaiotto:2008cd}：
$u$ 平面中各奇点周围 KS 因子的乘积给出恒等变换（更确切地说，是
红外电荷格的相应单值变换），但
$\zeta$ 平面中绕原点的所有 KS 因子的乘积则不然。
因此，没有任何形式线算子会带有
不适用于物理线算子生成函数的单值变换。

遗憾的是，正如上面的例子充分显示的，
直接与形式线算子打交道相当困难；因此我们愿意
提出：实际上有一个更强的猜想成立，使我们能与
更简单的对象打交道。我们令 $y=1$，并定义一个（可能更大的）
汇集 $\CL(c)$ 构成的空间，其中 $\CL(c)$ 定义在 $\Xi$ 的各腔中，
是 $x_\gamma$ 的正整系数线性组合，且在相邻的
腔之间 $\CL(c^\pm)$ 由适当的 $y=1$ KS 变换相联系。
我们不妨把这个空间的元素称为\emph{形式叶层}（formal lamination）。
「叶层」这个名字的由来将在第 10 节中变得清楚。% TODO 合稿改 \ref{sec:Examples}
形式叶层的空间容易研究，因为它带有自然的乘积，
且通常由有限个生成元生成。注意，
若把 (3.31) 中的乘积展开，展开式中只出现正的% TODO 合稿改 \eqref{eq:y-one-prod}
整数系数（对 $\gamma$ 的适当符号而言）。
于是正性与物理预期很好地相容。

当然，只要令 $y=1$，形式线算子就可以看成形式叶层。
不过并不显然的是：两个形式线算子的非交换乘积会产生形式线算子的线性组合
\begin{equation}
F_i \circ_y F_j = \sum_{k} c_{ij}^k(y) F_k
\end{equation}
另一方面，形式叶层的集合对乘法封闭则是显然的。在我们关注的 $A_1$ 例子中，形式叶层的空间
似乎与规范理论中\emph{物理}线算子的空间重合。构造形式叶层
比构造形式线算子容易得多！因此，如果形式
叶层的空间已经与形式线算子的空间重合，那将是一个非常有用的
结果。我们因此愿意提出：这是一个值得追问的问题。
